\section{Conclusion}
\label{sec:conclusion}

A transformer language model is iterated in three ways: over layers, over optimizer steps and over generated tokens. Each iteration has been studied with tools from dynamical systems, mostly by different communities and under different names. We gave the three iterations one formal frame, reviewed the results on each, and summarized them in a dictionary that records which concepts have been proved, measured or only conjectured on which clock. In experiments on open models we found that training makes each layer of a Pythia model more expansive while the coherence of the whole forward pass falls, and that gradient descent on a small transformer stays at a multi-mode edge of stability across two decades of step size. Repetition loops in GPT-2 and Pythia behave as wells whose depth grows while the model sits in them and then saturates. Whether a loop locks therefore depends on the length of the generation and on the size of the model, and coherent text lies in a narrow temperature window between repetition and incoherence. A small reasoning model has the same window, centered on its recommended temperature, and its low-temperature failure is recurrent rethinking, with no periodic loop.

A bridge between two fields can describe, explain and solve. This paper describes the three clocks in one vocabulary, shows the correspondences side by side, derives the reduced-order models behind them, and explains several behaviors that the machine learning literature had measured but not accounted for, from the Arrhenius escape from loops to the narrow window of coherent text. It does not yet solve a problem of language models that was unsolved before: our controllers for repetition find the right temperature but do not beat it. We think the third level is within reach, because many problems of language models have counterparts that dynamics has already solved, and the open problems of Section~\ref{sec:open} point to where.

Many cells of the dictionary are empty, and several of the open problems sit where two clocks meet, such as the effect of training on the depth dynamics, the feedback of generated text into training, and the prediction of training instabilities with the tools used for escape problems in mechanics. Each of these questions has a counterpart that nonlinear dynamics has studied for decades on pendulums, ships and oscillator networks, and each can be tested with the weights and outputs of open models on a laptop. We hope that the shared vocabulary of this paper helps researchers from both communities recognize the same phenomenon when they meet it, and solve it together.

\subsubsection*{Code and data}
Code to reproduce every figure and number, the raw results, and a maintained list of the papers in this review are available \codeurl. All experiments run on a CPU.

\subsubsection*{Broader impact statement}
This is a review with small-scale experiments on public models, and we do not foresee direct negative consequences. A better account of repetition loops and of training instabilities could reduce wasted compute in pretraining and make generation more reliable.

\subsubsection*{Use of AI tools}
\label{sec:ai}
An AI assistant (Claude, Anthropic) was used to search and verify the literature, draft and edit text, and write and run experiment code under the author's direction. All references were checked against arXiv or Crossref records, all numbers come from the released code, and the author is responsible for every claim in the paper.
